\setcounter{footnote}{23}

\printedpage{19}
% head: A. Grothendieck : Dérivateurs, Ch. XIII, Édition des Universités de Montpellier II et Paris VII
% foot: Version 25 janvier 2004
\msp{page 31}
as a \textit{right adjoint} (fully faithful) of the localization functor
\[
\pi_{\mathrm{d}}\cM_{\mathrm{c}} \longrightarrow (\pi_{\mathrm{d}}\cM_{\mathrm{c}})W_{\mathrm{c}}^{-1} \;\simeq\; \mathrm{Ho}_W,
\]
and $\mathrm{Ho}_W$ is obtained from $\pi_{\mathrm{d}}\cM_{\mathrm{c}}$ by a \textit{calculus of left fractions.} Dually, if we have right exactness, \textit{i.e.} $TF = F\cap W$ ($\Longleftrightarrow$ $TF = F\cap W_{\mathrm{univ}}$), then the relevant equivalence relation on $\mathrm{Hom}_{\cM}(X,Y)$, for $X$, $Y$ in $\cM_{\mathrm{cf}}$, is left Q-homotopy (implied by right Q-homotopy), and this time the natural arrow from $\pi_{\mathrm{g}}\cM_{\mathrm{cf}} = \pi\cM_{\mathrm{cf}} \simeq \mathrm{Ho}_W$ to $\pi_{\mathrm{d}}\cM_{\mathrm{f}}$ \texttt{[rather }$\pi_{\mathtt{g}}\cM_{\mathtt{f}}$\texttt{]},
\[
\underbrace{\pi\cM_{\mathrm{cf}}}_{\simeq\,\mathrm{Ho}_W} \longrightarrow \pi_{\mathrm{d}}\cM_{\mathrm{f}} \qquad \texttt{[rather }\pi_{\mathtt{g}}\cM_{\mathtt{f}}\texttt{]}
\]
is a \textit{left adjoint} (fully faithful) of the localization functor
\[
\pi_{\mathrm{d}}\cM_f \longrightarrow \cM_{\mathrm{f}}W_{\mathrm{f}}^{-1} \;\simeq\; \mathrm{Ho}_W \qquad \texttt{[rather }\pi_{\mathtt{g}}\cM_f \longrightarrow \cM_{\mathtt{f}}W_{\mathtt{f}}^{-1} \;\simeq\; \mathrm{Ho}_W\texttt{]}, % [sic] \cM_f (italic f, twice) -> \cM_{\mathrm{f}}, elsewhere the index f is upright
\]
and the localized category is made explicit by a calculus of right fractions.

In the biexact case (\textit{i.e.} that of genuine Quillen structures), the two Q-homotopy relations in $\cM_{\mathrm{cf}}$ (more generally in $\mathrm{Hom}(X,Y)$, if $X$ in $\cM_{\mathrm{c}}$ and $Y$ in $\cM_{\mathrm{f}}$)
\msp{page 32}
coincide, and even become particularly simple. (For if $f, g : X \rightrightarrows Y$, with $X$ cofibrant and $Y$ fibrant, then $f\sim g$ can be tested on \textit{any} ``cylinder object'' for $X$, or on \textit{any} ``path object'' for $Y$.)

In the case where the weak Quillen structure is left exact, \textit{i.e.} where the relevant equivalence relation in $\mathrm{Fl}\,\cM_{\mathrm{cf}}$ is right Q-homotopy, then the exactness axiom is satisfied for $W$-cartesian squares, but I doubt that the same holds for $W$-cocartesian squares (\footnote{It seems nevertheless that this is the case, according to \textsc{Baues}, cf. corresponding marginal annotation page 30.}). (Unless we make an additional hypothesis, e.g. right exactness of the Quillen structure.) If we have only exactness (but neither left nor right), it seems that we cannot infer from it the exactness either of $W$-cartesian squares or of $W$-cocartesian squares, hence not that of cosuspension sequences (loop sequences), nor of suspension sequences. We might therefore think that the exactness condition, all by itself, is not of much
\msp{page 33}
use for the homotopy formalism in the categories of fractions $\mathrm{Ho}_W$ etc. But we do nevertheless have this: let $\Psi = {}$\begin{tikzcd}[row sep=tiny, column sep=small, baseline=(\tikzcdmatrixname-2-1.base)]
 & 1 \\
0 \arrow[ru] \arrow[rd] & \\
 & 2
\end{tikzcd}, and suppose that $(\cM,W)$ \textit{and} $(\cM(\Psi),W(\Psi))$ \texttt{[where for every small category }$I$\texttt{, }$\cM(I)$\texttt{ denotes the category of presheaves on }$I$\texttt{ with values in }$\cM$\texttt{, and }$W(I)$\texttt{ is the subset of }\texttt{Fl}$(\cM(I))$\texttt{ consisting of morphisms}
\printedpage{20}
% head: A. Grothendieck : Dérivateurs, Ch. XIII, Édition des Universités de Montpellier II et Paris VII
% foot: Version 25 janvier 2004
\texttt{of presheaves that are in }$W$\texttt{ ``argument by argument'']} admit an exact weak Quillen structure. Then in $\cM$ the $W$-cartesian squares are $W$-exact, so we have in $\cM$ long exact cosuspension sequences (exact loop sequences) (\footnote{it seems that I botched it here, and I doubt that the result mentioned is correct.}). (The hypothesis on $(\cM(\Psi),W(\Psi))$ serves only to ensure that every arrow of $\mathrm{Ho}_{W(\Psi)} = \cM(\Psi)W(\Psi)^{-1}$ is isomorphic to an arrow coming from an arrow of $\cM(\Psi)$ itself. But this is an apparently very delicate property, and I know of no other way to obtain it than via the existence of an exact weak Quillen structure.) Dually, if $(\cM,W)$ and $(\cM(\Phi),W(\Phi))$ admit exact weak Quillen structures, then the $W$-cocartesian squares in $\cM$ are $W$-exact, and we find (if $\cM$ is pointed) long exact suspension sequences. % [sic] W-exactes -> W-exacts, agreement with "carrés" (cf. "W-exacts" below)

\msp{page 34}
This shows that if, for a given category of models $(\cM,W)$, we are interested in ``doing homotopical algebra'' for all the pairs $(\cM(I),W(I))$ at once, with $I$ in Cat, it is enough (if we want to have a good theory of $W$-homotopic fibers and cofibers in these categories, including the property of $W$-exactness for $W$-cartesian squares and for $W$-cocartesian squares) to have for each $(\cM(I),W(I))$ an exact weak Quillen structure, nothing more.

This therefore shows us the interest of developing criteria, for a $(\cM,W)$, ensuring that this is indeed so. On this subject, I note that I know how to check the exactness of a Quillen structure \textit{only if} I have left exactness, or right exactness. This means that I must in fact bring out conditions on $(\cM,W)$ ensuring (say) that the $(\cM(I),W(I))$ admit left exact weak Quillen structures.

\msp{page 35}
The weak Quillen structures for the $(\cM(I),W(I))$ will be deduced from a weak Quillen structure for $(\cM,W)$, in the following way: they will all be \textit{``right closed'',} \textit{i.e.}
\[
F(I) \;=\; (TC(I))_*, \qquad TF(I) \;=\; (C(I))_*, \tag*{(1)}
\]
in particular (for $I=e$)
\[
F \;=\; (TC)_*, \qquad TF \;=\; C_*,
\]
which already implies (since $C\subset \mathrm{Cof}_W$, $TC\subset W^{\mathrm{univ}}$)
\[
(\mathrm{Cof}_W)_* \;\subset\; W, \qquad (W^{\mathrm{univ}})_* \;\subset\; \mathrm{Fib}_W \tag*{(2)}
\]
(whence $(\mathrm{Cof}_W)_*\subset W_{\mathrm{univ}}$). By virtue of (1), to define the weak Quillen structures for the $(\cM(I),W(I))$, it suffices to define $C(I)$, $TC(I)$, and even, since we want left exactness, \textit{i.e.}
\[
TC(I) \;=\; C(I)\cap W(I) \tag*{(3)}
\]
\printedpage{21}
% head: A. Grothendieck : Dérivateurs, Ch. XIII, Édition des Universités de Montpellier II et Paris VII
% foot: Version 25 janvier 2004
(in particular
\[
TC \;=\; C\cap W \quad ), \tag*{(4)}
\]
it remains to define $C$. We shall define $C(I)$ ``argument by argument'', \textit{i.e.} for $f\in\mathrm{Fl}(\cM(I))$
\[
f\in C(I) \qquad\Longleftrightarrow\qquad \alpha_s^*(f)\in C \ \text{ for every } s\in\mathrm{Ob}\,I, \tag*{(5)}
\]
where
\[
\alpha_s : e\longrightarrow I \qquad \texttt{[is]}\ \text{defined by } \alpha_s(e)=s.
\]
It follows moreover at once, by (3),
\[
f\in TC(I) \qquad\Longleftrightarrow\qquad \alpha_s^*(f)\in TC \ \text{ for every } s\in\mathrm{Ob}\,I. \tag*{(6)}
\]

\msp{page 36}
We therefore suppose that $(\cM,W)$ admits a left exact and right closed weak Quillen structure $(C,TC,F,TF)$, defined in terms of
\[
C\;\subset\;\mathrm{Fl}(\cM),
\]
by
\[
TC \;=\; C\cap W, \qquad F \;=\; (TC)_*, \qquad TF = C_*.
\]
For this to really give a weak Quillen structure (necessarily left exact and right closed), we simply have to check the following conditions.

\begin{itemize}[leftmargin=3em]
\item[\circled{i}] Factorization for $(C, TF=C_*)$ and $(TC=C\cap W, F=(TC)_*)$ – we shall come back to this point.
\item[\circled{ii}] Stability of $C$ under composition and cobase change, and $C$ contains the isomorphisms (from which the similar stability of $TC = C\cap W$ will follow, taking into account that $C\subset\mathrm{Cof}_W$) (NB The appropriate stabilities for $F$, $TF$ are trivial by definition $F=(TC)_*$, $TF=C_*$.)
\item[\circled{iii}] $C\subset\mathrm{Cof}_W$, $\underbrace{C_*}_{=\,TF}\subset W$, $(C\cap W)_*\subset\mathrm{Fib}_W$.
\end{itemize}

To check, similarly, that $(C(I),TC(I),F(I),TF(I))$ is a weak Quillen structure for $W(I)$, we must check the conditions similar to \circled{i}\,\circled{ii}\,\circled{iii} . Now (\textit{ii}) for $C(I)$ follows at once from condition (\textit{ii}) for $C$. It remains to examine (\textit{iii}) and (\textit{i}).

\printedpage{22}
% head: A. Grothendieck : Dérivateurs, Ch. XIII, Édition des Universités de Montpellier II et Paris VII
% foot: Version 25 janvier 2004
\msp{page 37}
The inclusions of (\textit{iii}).

The inclusion
\[
C\;\subset\;\mathrm{Cof}_W
\]
implies that every $f\in C(I)$ is in $\mathrm{Cof}_W$ argument by argument. This implies that $f$ is in $\mathrm{Cof}_{W(I)}$, as we see at once from the definition – I make it explicit in the

\textbf{Lemma 1:} \textit{Let $f : X_\bullet\longrightarrow Y_\bullet$ be an arrow of $\cM(I)$. For $f$ to be in} $\mathrm{Cof}_{W(I)}$ (resp. \textit{in} $\mathrm{Fib}_{W(I)}$), \textit{it} suffices \textit{that for every $s$ in} $\mathrm{Ob}\,I$, \textit{$\alpha_s^*(f)$ be in} $\mathrm{Cof}_W$ (resp. \textit{in} $\mathrm{Fib}_W$)\textit{. Likewise (for the record) for a condition for $f\in W(I)$ – there it is necessary and sufficient that $\alpha_s^*(f)\in W\ \forall\, s\in\mathrm{Ob}\,I$.}

We shall apply this same lemma to establish the inclusions
\[
C(I)_* \;\subset\; W(I), \qquad (C(I)\cap W(I))_* \;\subset\; \mathrm{Fib}_{W(I)}.
\]
So it suffices to see that
\[
\begin{aligned}
\alpha_s^*(C(I)_*) \;&\subset\; C_* &&\qquad (\subset\; W)\\
\alpha_s^*((\underbrace{C(I)\cap W(I)}_{(C\cap W)(I)})_*) \;&\subset\; (C\cap W)_* &&\qquad (\subset\; \mathrm{Fib}_W),
\end{aligned}
\]
and more generally
\[
\alpha_s^*(\Phi(I)_*) \;\subset\; \Phi_* \tag*{(7)}
\]
for any ( ?) subset $\Phi$ of $\mathrm{Fl}(\cM)$, thus defining $\Phi(I)\subset\mathrm{Fl}(\cM(I))$ term by term.

\msp{page 38}
Making this explicit, we find the condition
\[
\alpha_{s\,!}(\Phi)\subset\Phi(I) \qquad \text{for every } s\in\mathrm{Ob}\,I,
\]
for it to work, \textit{i.e.} we want that
\[
i\in\Phi \qquad\Longrightarrow\qquad \alpha_{s\,!}(i)\in\Phi(I), \quad \textit{i.e.}\ \ \alpha_t^*(\alpha_{s\,!}(i))\in\Phi\ \ \forall\, t\in\mathrm{Ob}\,I.
\]
Now the functor $\alpha_t^*\alpha_{s\,!}$ is made explicit as
\[
A \longmapsto \underbrace{\mathrm{Hom}_I(t,s)\times A}_{\substack{\text{sum of } \mathrm{Hom}_I(t,s)\\ \text{copies of } A \text{ in } \cM}},
\]
so we must suppose $\Phi$ \textit{stable under direct sums} (of the type considered $\mathrm{Hom}_I(t,s)\times A$). This is indeed the case for $\Phi = C$ and $\Phi = C\cap W$ and for direct sums that are \textit{finite,} thanks to (\textit{ii}). We thus find the

\printedpage{23}
% head: A. Grothendieck : Dérivateurs, Ch. XIII, Édition des Universités de Montpellier II et Paris VII
% foot: Version 25 janvier 2004
\textbf{Lemma 2.}

\begin{itemize}[leftmargin=3em]
\item[a)] \textit{Let $I$ in} Cat \textit{be such that direct sums of type} $\mathrm{Hom}_I(t,s)$ \textit{exist in $\cM$. (E.g. we suppose that $I$ is} finite, \textit{more generally that the} $\mathrm{Hom}_I(t,s)$ \textit{are finite sets; or that ``small'' direct sums exist in $\cM$.) Let $\Phi$ be a subset of} $\mathrm{Fl}(\cM)$\textit{, stable under this type of direct sums, and let} $\Phi(I)\subset\mathrm{Fl}(\cM(I))$ \textit{be the corresponding subset. Then for every $s\in\mathrm{Ob}\,I$, we have}

\msp{page 39}
\[
\alpha_s^*(\Phi(I)_*) \;\subset\; \Phi_*.
\]
\item[b)] \textit{Suppose that the} $\mathrm{Hom}_I(t,s)$ \textit{are finite (e.g. $I$ finite). Then the hypothesis of} a) \textit{applies for $\Phi=C$ and for $\Phi=C{\cap}W$, so condition} (iii) \textit{p.} 36 \textit{is satisfied for the structure $(W(I),C(I),TC(I),F(I),TF(I))$ in $\cM(I)$. Same conclusion without a finiteness hypothesis on the} $\mathrm{Hom}_I(t,s)$\textit{, when we suppose $\cM$ stable under ``small'' direct sums, and $C$ and $W$ also stable under such direct sums.}
\end{itemize}

\textbf{NB} Given the scholie above, it is perfectly reasonable to restrict ourselves if need be to categories $I$ that are \textit{finite} – given that the category of these $I$ is stable under the products $I\longmapsto I\times\Psi$, $I\longmapsto I\times\Phi$.

It remains, to be happy, to check the factorization condition \circled{i} for $(C,TF)$, $(TC,F)$, but also and above all for the $(C(I),TF(I),TC(I),F(I))$. This means that even starting from a closed Quillen structure (land of Cockaigne)
\msp{page 40}
does not get us very far all by itself, for lack of knowing that the factorization axiom is preserved, in passing from $(\cM,W,\ldots)$ to $(\cM(I),W(I),\ldots)$.

It is here that the factorization theorem of no.~1 should prove providential. But to be able to apply it to the pairs $(C(I),C(I)_*)$, $(TC(I),TC(I)_*)$, we need (in addition to the accessibility of $\cM$, which is a quite innocuous condition) very strong stability conditions for $C(I)$, $TC(I)$ (\footnote{Plus the \textit{small} systems of generators for $C$, $TC$.}). To tell the truth, there remains the stability of $C$ and of $C\cap W$ under ordinal direct limits (which will imply the same stability for the $C(I)$, $(C{\cap}W)(I) = C(I){\cap}W(I)$).

Truth to tell, I see little chance of finding a weak Quillen structure $(C,TC,F,TF)$ associated with $W$, and satisfying all the stated conditions, including that of stability of $C$ and of $C\cap W$ under ordinal $\varinjlim$, without supposing that $W$ itself is already stable under (small) filtering direct limits. % [sic] associé -> associée, agreement with "structure"

\msp{page 41}
We can think of using a \textit{functorial} factorization of the arrows in $\cM$ as
\[
\begin{aligned}
f = pi \quad \text{with}\quad & i\in TC, && p\in F\\
\text{or}\quad & i\in C, && p\in TF,
\end{aligned}
\]